% Isolated proof.  This file is not a live-book input.
% It has not been submitted to Rethlas.

Choose positive numbers \(\delta_n\) such that
\(\sum_{n\geq1}\delta_n<\infty\), and put
\[
                         c_n=\sum_{m\geq n}\delta_m.
                         \tag{3}
\]
We construct indices \(k_n\) and put \(w_n=v_{n,k_n}\).  Choose \(k_1\)
arbitrarily.  Suppose \(w_n\) has been chosen.  We claim that one can choose
\(k_{n+1}\) so that, on \(K\),
\begin{align}
 v_{n+1,k_{n+1}}&<w_n+\delta_n,                         \tag{4}\\
 v_{n+1,k_{n+1}}&<v_{r,s}+\delta_n
                    &&(1\leq r,s\leq n).                \tag{5}
\end{align}

Indeed, all functions occurring in these finitely many inequalities are
continuous and are defined on a common open neighbourhood of \(K\).  At
each \(y\in K\),
\[
 u_{n+1}(y)\leq u_n(y)\leq w_n(y)
 \quad\hbox{and}\quad
 u_{n+1}(y)\leq u_r(y)\leq v_{r,s}(y)\quad(1\leq r,s\leq n).
\]
Since \(v_{n+1,k}(y)\downarrow u_{n+1}(y)\), some \(k\), depending at first
on \(y\), makes all the strict inequalities in (4)--(5) true at \(y\).
For each \(k\), the locus on which all these inequalities hold is open, and
these loci increase with \(k\) and cover \(K\).  Compactness of \(K\) gives
one \(k_{n+1}\) that works on all of \(K\).  This proves the claim.

On \(K^\circ\), set
\[
                              g_n=w_n+c_n.               \tag{6}
\]
Restriction preserves smoothness and positivity of
\(\mathrm d'\mathrm d''\), and addition of a constant does as well.  By
(3)--(4),
\[
 g_{n+1}<w_n+\delta_n+c_{n+1}=w_n+c_n=g_n
\]
on \(K^\circ\).  Thus \((g_n)\) is a decreasing sequence of smooth psh
functions on one fixed neighbourhood.

It remains to identify its limit; this is the point not supplied by merely
imposing (4).  Because every term of a decreasing row dominates its row
limit,
\[
                       g_n\geq w_n\geq u_n\geq u.        \tag{7}
\]
Let \(g=\inf_n g_n\).  Fix \(r,s\geq1\).  For every \(N\) sufficiently
large, (5), used when \(w_N\) was selected, gives
\[
                      w_N<v_{r,s}+\delta_{N-1}.
\]
Since \(c_N\to0\) and \(\delta_{N-1}\to0\), passage to the limit yields
\(g\leq v_{r,s}\) on \(K^\circ\).  First let \(s\to\infty\), and then let
\(r\to\infty\).  The two row limits give
\[
                         g\leq u_r\quad\hbox{for every }r,
                         \qquad g\leq u.                 \tag{8}
\]
Together, (7) and (8) prove \(g=u\), hence (1).

Under the final hypothesis in the statement, the argument applies at every
point of \(U\).  Therefore a decreasing limit of sections of
\(\operatorname{Psh}_{X,1}(U)\) is again a section of
\(\operatorname{Psh}_{X,1}(U)\).  The latter sheaf contains
\(\operatorname{Psh}_{X,0}\) and is now stable under the operation in
Chambert--Loir--Ducros, Definition~8.1.1.  Its defining minimality gives
\(\operatorname{Psh}_X\subseteq\operatorname{Psh}_{X,1}\); the reverse
inclusion is built into the transfinite construction of
Chambert--Loir--Ducros, Lemma~8.1.4.  This proves (2).

